\section{Introduction}

The Fractions Skill Score (FSS) is a neighbourhood-based verification metric
for spatial forecasts, introduced by \citet{roberts2008}. Rather than comparing
forecast and observation values point by point, it compares the \emph{fraction}
of grid points that exceed a given threshold within a sliding window
(neighbourhood), as a function of both the threshold and the window
(neighbourhood) size. This avoids the \emph{double penalty} that afflicts
point-wise scores when a feature is forecast with a small spatial displacement
and reveals the spatial scale at which a forecast attains useful skill. It does
so by treating the spatial distribution of events within a window
probabilistically: a perfect forecast is one with the same event frequency as
the observation within the window, regardless of the exact placement.

For a single threshold and window size the score is \citep{roberts2008}
%
\begin{equation}
  \mathrm{FSS} = 1 -
  \frac{\dfrac{1}{N}\displaystyle\sum_{i=1}^{N}\bigl(p_f - p_o\bigr)^2}
       {\dfrac{1}{N}\displaystyle\sum_{i=1}^{N}p_f^{\,2}
        + \dfrac{1}{N}\displaystyle\sum_{i=1}^{N}p_o^{\,2}},
  \label{eq:fss}
\end{equation}
%
where $N$ is the number of windows in the domain (one centred on each grid
point), $p_f$ is the forecast fraction and $p_o$ the observed fraction of the
window. $\mathrm{FSS}=1$ denotes a perfect forecast and $\mathrm{FSS}=0$ no
skill.

The cost of the FSS comes from evaluating the windowed fractions $p_f$, $p_o$
for many window sizes and thresholds. \citet{faggian2015} showed that this
windowed summation can be computed very efficiently using a \emph{summed area
table} (also called an integral image): the table is built once per threshold
with two cumulative sums, after which the sum over any rectangular window is
obtained from just four array lookups, independent of the window size. This is
the algorithm implemented in \texttt{fss\_SAT.py}.

\subsection{Scope of this document}

\texttt{fss\_SAT.py} started from the reference Python implementation in the
appendix of \citet{faggian2015}. The purpose of this document is to record the
\emph{adaptations and extensions} that were made on top of that baseline. It is
organised as follows:

\begin{itemize}
  \item Section~\ref{sec:baseline} recaps the SAT algorithm of
        \citet{faggian2015} and its reference implementation, establishing the
        notation and the starting point for the changes.
  \item Section~\ref{sec:adaptations} documents each adaptation made in
        \texttt{fss\_SAT.py}: the performance rewrites that preserve the
        original result bit-for-bit, the feature extensions (threshold modes,
        percentile thresholds, ensemble FSS, parallelisation), and the
        substantive algorithmic additions (missing-data handling and the
        continuous/weighted FSS).
\end{itemize}

A companion FFT-based implementation (\texttt{fss\_FFT.py}) and an OpenCL
implementation (\texttt{fss\_OCL.py}) exist in the same repository and compute
the same score; they are out of scope here and documented separately.
